\SourceLeaf{390}%
\noindent hervorgeht:
\[
Q'=F(t,u)=P,
\]
$A,B$ und $C$ Functionen von $P$ sein müssen.

[Aus der Gleichung $qt+Qu=1$ findet man durch Differentiation nach $t$ und $u$:
\[
\begin{aligned}
q+\left(t+\frac{dQ}{dq}u\right)q'&=0\\
Q+\left(t+\frac{dQ}{dq}u\right)q''&=0\\
\left(2+\frac{ddQ}{dq^2}q'u\right)q'+\left(t+\frac{dQ}{dq}u\right)q^{(2)}&=0\\
\frac{dQ}{dq}q'+q''+\frac{ddQ}{dq^2}u q'q''+\left(t+\frac{dQ}{dq}u\right)q^{(3)}&=0\\
\left(2\frac{dQ}{dq}+\frac{ddQ}{dq^2}q''u\right)q''+\left(t+\frac{dQ}{dq}u\right)q^{(4)}&=0
\end{aligned}
\]
und hieraus folgen die Identitäten:
\[
q''q''q^{(2)}-2q''q'q^{(3)}+q'q'q^{(4)}=0
\]
und
\[
P''P''P^{(2)}-2P''P'P^{(3)}+P'P'P^{(4)}=0.]
\]
\par\indent Diess gibt
\[
A''A''A^{(2)}-2A'A''A^{(3)}+A'A'A^{(4)}=0;
\]
allein man sieht noch nicht, wie diese Gleichung lediglich vermittelst des Obigen aus $m=1$ abgeleitet werden könnte.
\Needspace{10\baselineskip}\begin{center}[5.]\end{center}
\noindent\hspace*{1.2em} 15. Dec. Diess ist nunmehro abgemacht, obwohl man dadurch für den Fall, wo $m$ nicht constant ist, nichts zu gewinnen vermag.

Es wird nemlich [wenn $m=1$ ist]
\[
\BareRoot{\{(A'dt+A''du)^2+(B'dt+B''du)^2+(C'dt+C''du)^2\}}
\]
ein Binomium von der Form $p'dt+p''du$ und ein vollständiges Differential, welches wir $=dq$ setzen. Diess scheint am einfachsten aus $\Delta=0$ und $\FourRingCross$ zu folgen.
